% =============================================================================
% 7. POLICY RECOMMENDATIONS  (Reviewer comment #2)
% =============================================================================
\section{Policy recommendations}
\label{sec:policy}

Building on the empirical case and the managerial implications, three
policy recommendations are formulated for civil aviation authorities and
industry bodies. Each recommendation targets a specific actor and a
specific instrument, and each is calibrated to the current state of \ICAO{}
Annex 19 and adjacent regulatory frameworks.

\subsection{Recommendation 1 -- Embed training-effectiveness evaluation in SMS acceptance criteria}
\label{subsec:pol-rec1}

Civil aviation authorities (\FAA{}, \EASA{}, \ANAC{} and equivalents) are
encouraged to require, as part of the periodic acceptance of \SMS{} manuals
submitted by airlines, airports and \MRO{} facilities, evidence of training
evaluation across all four Kirkpatrick levels. The current text of
\ICAO{} Annex 19 leaves the depth of evaluation to the discretion of the
service provider, which produces wide variance in practice
\citep{Nakamura2024,ORconnor2024}. A targeted clarification --- requiring
that Level~3 and Level~4 metrics be reported alongside Level~1 and Level~2
in the annual \SMS{} performance dashboard --- would close the variance
without imposing a prescriptive methodology.

\subsection{Recommendation 2 -- Industry guidance on training evaluation for MRO}
\label{subsec:pol-rec2}

The IATA and SAE are encouraged to develop a joint
guidance document on training evaluation for \MRO{}, modelled on the
existing safety-performance indicator (SPI) guidance. The
document should specify minimum data sources, recommended metrics by
Kirkpatrick level, and a maturity model that organisations can use to
self-assess and to communicate with auditors. A worked example based on
\AS{} clauses (as in this study) would accelerate adoption.

\subsection{Recommendation 3 -- Mandate annual reporting of training-effectiveness metrics as safety indicators}
\label{subsec:pol-rec3}

Aviation safety regulators are encouraged to mandate annual reporting of
training-effectiveness metrics (Level~3 and Level~4) as part of the safety
performance indicator portfolio that organisations submit. Such reporting
need not be public to be effective --- regulator-only access already
delivers most of the accountability benefit \citep{FAA2024}. Aggregated and
anonymised disclosure to the industry would further support
benchmarking and the diffusion of best practice.

\paragraph{Coupling with sustainability disclosure.}
Where airlines or \MRO{} facilities are subject to \ESG{} reporting
regimes (for example, the ESRS in the European Union), the
training-effectiveness metrics can be cross-referenced from the safety
performance dashboard to the human-capital and cleaner-production sections
of the sustainability report. This cross-referencing reduces reporting
duplication and reinforces the safety--sustainability nexus that motivates
the present study.
